
We demonstrated in proof-of-concept that verifier feedback provides a measurable semantic gradient for LLM specification synthesis. By decomposing feedback into three informational dimensions---Witness Instantiation, Logical Structure, Dependency Preservation---and abstracting via an 8-primitive taxonomy, we enable systematic iterative refinement with cross-verifier portability.

Our contributions extend verification-in-loop from empirical observation to quantitative analysis through information-theoretic framework. The quantified information gradient ($\beta =12.49, R^{2}=0.89)$ provides basis for feedback design; cross-verifier retention (84.9\%) demonstrates semantic normalization preserves utility; compute-matched control isolates feedback as causal mechanism.

The bottleneck shifts from ''LLMs cannot do formal reasoning'' to ''we must design feedback as first-class learning signal.'' Research frontiers emerge: (1) Learned semantic normalization---replace hand-crafted taxonomy with learned abstractions, (2) Verifier-LLM co-design---optimize proof obligation structure for LLM interpretability, (3) Probabilistic correctness---combine formal verification with learned confidence estimation.

Viewing verification-in-loop through information theory opens research directions where verification and learning are complements rather than opposites.
